\documentclass{article}
\usepackage{amsmath,amssymb,amsthm}
\newtheorem{lemma}{Lemma}\newtheorem{theorem}{Theorem}\newtheorem{remark}{Remark}
\newcommand{\N}{\mathcal N}
\begin{document}
% Standalone analytic fragment. No numerical coverage is asserted here.
\begin{lemma}[Exact pricing on physical source polygons]
\label{lem:physical-source-polygon-price}
Fix an amplitude interval $[u_-,u_+]$, a source interval $[r_-,r_+]$,
and a cumulant interval. Suppose that, throughout the cumulant interval
and for every allowed shape, proved scalar inequalities imply
\[
 a_j u+r\le C_j,\qquad a_j>0\quad(1\le j\le m).
\]
Let $K$ be the intersection of these half-planes with the amplitude--source
rectangle, after any separately proved coordinate restrictions.
For rational inputs its vertices and the maximum on $K$ of
\[
 F(u,r)=u^2+L_u u+g r^2+L_r r+C
\]
are computable using rational arithmetic. For a nonempty $K$, it suffices
to check its vertices and, on each edge, the stationary point in the edge
interior when the restricted quadratic is concave. This includes segments
and singletons. An empty intersection contains no physical source.

For a native certificate interpolated in $r$, first split $K$ at every
source knot. On each resulting piece apply this rule to each of the four
cumulant--shape corner bounds, and take the maximum of the resulting
bounds. A fixed nonnegative contact multiplier may be included by changing
only $L_u,L_r,C$.
\end{lemma}
\begin{proof}
Every defining half-plane contains every physical source in the nominal
rectangle, so intersection and coordinate clipping preserve all physical
points. Rational line intersections give all vertices, retaining points
on the boundaries. If a full-dimensional polygon had a maximizing point
in its interior, sufficiently small horizontal displacements would stay
in the polygon. Their two objective values have average
$F(u,r)+\delta^2$, contradicting maximality. Thus its maximum lies on its
boundary. On the edge $(u,r)=(u_0,r_0)+t(\Delta u,\Delta r)$,
$0\le t\le1$, the quadratic coefficient is
$(\Delta u)^2+g(\Delta r)^2$. Endpoints suffice when this coefficient is
nonnegative; otherwise also check the unique stationary point if it lies
in $(0,1)$. All these operations are rational. The same calculation covers
the lower-dimensional cases.

A native source chord is valid only between successive saved knots.
Splitting at every knot therefore retains its validity on each closed
piece; shared boundaries have both valid endpoint bounds. For each fixed
cumulant--shape corner the chord changes only the linear and constant
coefficients in $r$. The already proved corner inequalities and positive
interpolation uniformly bound the whole cumulant--shape cell. Taking the
maximum over the four corner prices and all source pieces gives a uniform
bound. The scalar row chosen at a source endpoint is fixed throughout its
cumulant--shape cell; different endpoints may use different globally
valid scalar rows.

For a contact with parameters $a,B$, put $d^2=6B^2$ and let $A$ be its
uniform scalar price. The nonnegative multiplier $\lambda$ adds
\[
 -\lambda a u-\lambda r+
 \lambda\left(A+\frac{\beta d}{b}\right)
 +\frac{\lambda^2d^2}{4b}.
\]
Here $\beta$ is the normalized native Hermite coefficient; for a raw
Hermite R3 profile the product $\beta d$ is $6\beta_{\rm raw}B$.
The coefficient of $u^2$ remains one, so the same boundary computation
applies. The multiplier and scalar contact are fixed on the whole
cell. Selecting any one of several proved upper bounds on each cell
preserves validity.
\end{proof}

\begin{remark}[Uniform source planes]
For the accepted even sextic scalar rows, the moment price is affine in
shape and convex in cumulant: the sextic coefficient is nonnegative and
the coefficient multiplying $\kappa^{3/2}$ is positive throughout the
shape interval. Hence the largest of the four endpoint prices is a
uniform $C_j$. A single fixed scalar row is used for each plane over the
whole cumulant interval. Taking the smallest uniform price among several
such rows is valid; taking separate minima at the four corners would
require a separate justification.
\end{remark}
\begin{theorem}[Complete preferred scalar certificate]
\label{thm:final-preferred-polygon}
Set $\widehat b=713/1000$. Then
\[
 \N(\widehat b)\le\widehat B:=\frac{564479481}{10^9}=.564479481.
\]
Consequently the explicit reduced rational matrix sequence has actual limit
\[
 \lim_{n\to\infty}G(H_n(\widehat b))
 =\frac1{\N(\widehat b)}
 \ge\frac{10^9}{564479481}
 =1.7715435789241736\ldots .
\]
This rational number is an unconditional lower bound for $G_{\mathbb R}$.
\end{theorem}
\begin{proof}[Certificate proof]
First transfer the complete outer branches of
Theorem~\ref{thm:same-b7171-regions} from $\bar b=7171/10000$ to
$\widehat b$. The physical amplitude, source, cumulant and shape variables
do not depend on the operator coefficient. All ten native profiles used
by those branches satisfy the exact normalized longitudinal budget
$\beta^2\le1/16$; the sole R3 replacement satisfies the corresponding
raw budget $6\beta_{\rm raw}^2\le1/16$. The smaller-coefficient spectral
gates remain valid. The price change is
\[
 (\widehat b-\bar b)r^2+
 \beta_{\rm norm}^2(1/\widehat b-1/\bar b).
\]
Dropping its nonpositive first term therefore bounds both complete outer
branches by
\[
 \frac{56338}{10^5}+\frac1{16}
   \left(\frac1{\widehat b}-\frac1{\bar b}\right)
 =\frac{144154052987}{255646150000}
 <\widehat B.
\]
The elementary low-amplitude branch improves directly when $b$ decreases.
A separate independent checker verifies the entire selected outer cover,
its $45025$ native cells, $41$ R3 cells, four earlier transported frames,
profile budgets, source hashes and spectral gates.

On the remaining amplitude interval $[1381/2000,31/40]$, the frozen
cover has $21$ amplitude bands and $45$ cumulant bands. The $945$ groups
partition $\kappa\in[0,90]$ and cover every physical source from zero
to its outward $R_{.847}(u_-)$ bound. The native shape $s\in[0,1]$
corresponds exactly to $z=2/5+(3/5)s$ for R3 prices. Independent directed
checks verify all $21$ original source-grid upper endpoints and all
$250$ inverse source clips.

Fourteen accepted whole-line R3 profiles give $630$ cumulant-band
constants. A separate replay checks all $2520$ cumulant--shape corner
prices, compact and tail reserve bindings, and the six spectral matrix
gates at the common coefficient $\widehat b$. The seven native endpoint
tables have $9360$ independently replayed baseline prices. Their exact
transfer from $\bar b$ adds
$\beta^2(1/\widehat b-1/\bar b)$ to each endpoint price and
$\widehat b-\bar b$ to its source quadratic coefficient.

Each cumulant band has ten fixed, proved uniform positive-head source
planes. All $450$ planes and their $1800$ moment corner prices are
independently verified. Intersecting these planes and the other proved
source restrictions retains all physical points. Of the original
$187335$ cells, precisely $1574$ are empty and $185761$ retain a selected
upper price. Among the latter, $8612$ use complete recorded polygons;
the remaining cells use valid enclosing rectangles.

The producer-free exact replay checks every retained polygon by feasible
vertices and original-constraint facet identities, splits each selected
native chord at every saved source knot, and includes every concave
quadratic edge vertex. It reconstructs all selected rectangular,
polygonal and fixed-contact prices and accounts for every original cell
exactly once. All recorded caps agree, and their maximum is
\[
 \frac{282239739989}{5\cdot10^{11}}<\widehat B.
\]
The frozen source is
\texttt{global\_\allowbreak middle\_\allowbreak b713\_\allowbreak finekappa46\_\allowbreak fourteenbank\_\allowbreak polygon\_\allowbreak envelope.json.gz};
the exhaustive independent reconstruction is
\texttt{final46\_\allowbreak b713\_\allowbreak polygon\_\allowbreak envelope\_\allowbreak audit.json}.
Thus the middle and outer branches bound every odd sign pairing in
every finite Gaussian dimension. Replacing one sign by its negative
gives the absolute norm bound. The established Gaussian-to-matrix
convergence and exact ratio-preserving reductions give the stated
actual matrix limit and its lower bound.
\end{proof}
\end{document}